% TOP_PROBLEM: {"problemId": "problem.grothendieck-period-conjecture", "canonicalTitle": "Grothendieck Period Conjecture", "releaseRank": 49, "primaryDomain": "number_theory_arithmetic_geometry", "primaryDomainLabel": "Number theory & arithmetic geometry", "displayStatus": "open", "releaseStatus": "open"}
% !TeX program = lualatex
% Definition and sources only; this file is not an LLM solution attempt.
\documentclass[11pt]{article}
\usepackage[margin=25mm]{geometry}
\usepackage{fontspec}
\setmainfont{Latin Modern Roman}
\usepackage{amsmath,amssymb,mathtools}
\usepackage{unicode-math}
\setmathfont{Latin Modern Math}
\usepackage{xurl,hyperref}
\hypersetup{colorlinks=true,urlcolor=blue}
\setcounter{secnumdepth}{0}
\setlength{\parindent}{0pt}
\setlength{\parskip}{5pt}
\setlength{\emergencystretch}{3em}
\begin{document}
\section*{\texorpdfstring{Grothendieck Period Conjecture}{Grothendieck Period Conjecture}}\label{49}
\subsection{Definitions and mathematical statement}
For a Nori motive $M/{\mathbb{Q}}$, its Betti and de Rham realizations are finite-dimensional rational vector spaces compared over ${\mathbb{C}}$. Their comparison matrix, in rational bases, supplies the periods of $M$. Its motivic Galois group $G_{\rm mot}(M)$ is the affine algebraic group of tensor automorphisms of the fiber functor on the tensor category generated by $M$. The conjecture is
\[
 \operatorname{trdeg}_{{\mathbb{Q}}}{\mathbb{Q}}(\text{entries of the period matrix of }M)
 =\dim G_{\rm mot}(M).
\]
Changes of rational bases do not change the generated period field. The Nori-motive convention and the category generated by $M$ are those of the source.
\par\smallskip\textit{Formulation and concept references:} \href{https://www.cambridge.org/core/journals/forum-of-mathematics-sigma/article/functional-transcendence-of-periods-and-the-geometric-andregrothendieck-period-conjecture/6B3E39BB9BC68DC2BC6E793D15B0735B}{[S1]}.

\subsection{Short English statement}
Are all algebraic relations among a motive's periods exactly those forced by its motivic symmetries?
\subsection{Sources}
\begin{itemize}
\item[S1]Formal statement, Conjecture 1.1.
\url{https://www.cambridge.org/core/journals/forum-of-mathematics-sigma/article/functional-transcendence-of-periods-and-the-geometric-andregrothendieck-period-conjecture/6B3E39BB9BC68DC2BC6E793D15B0735B}.
\textit{Formulation source.}
\item[S2]Functional Transcendence of Periods and the Geometric André--Grothendieck Period Conjecture.
\url{https://arxiv.org/html/2208.05182v2}.
\textit{Further reference; not a proof endorsement.}
\item[S3]Grothendieck's period conjecture for Kummer surfaces of self-product CM type.
\url{https://arxiv.org/abs/2303.05030}.
\textit{Further reference; not a proof endorsement.}
\end{itemize}
\end{document}
